\subsection{Selling part of the hedged inventory}
\label{sec:hedged-inventory-subsequent}

The purchase has absorbed the entire hedge reserve, but some goods may remain
on the shelf when the company next reports. Their cost stays in inventory until
sale, subject to a write-down if net realisable value falls below that cost.
The sale does not release another portion of the hedge reserve: that reserve
was already removed when all the goods were acquired.

Take the preceding purchase of 130 units with adjusted cost 10800, and suppose
the units are ordinarily interchangeable, with the same nature and use. There
is no opening stock or additional purchase. We apply weighted-average cost,
one of the formulas permitted by IAS~2.25 and explained in IAS~2.27. The
original hedge designation layers describe when the exposure entered the
hedge; they do not identify separate physical inventory lots. For this one
uniform batch, the unit cost is
\[
 c=\frac{10800}{130}=\frac{1080}{13}.
\]
Selling 52 units for 5800 therefore recognizes cost $52c=4320$ and leaves
$78c=6480$ in inventory. In the positive-hedge case, cumulative pre-tax profit
is $90+5800-4320=1570$. The earlier 90 of hedge ineffectiveness remains in
profit; the adjustment of purchase cost has already done its separate work.

Now suppose reliable evidence supports net realisable value of 70 per unit
for the 78 units still held. This is a supplied estimate after completion and
selling costs, with the same selling circumstances for every unit. Inventory
falls from 6480 to $78\times 70=5460$, and the company recognizes a write-down
of 1020. IAS~2.28--31 requires the estimate to reflect the goods and their
purpose; different contracts or selling circumstances would require a more
detailed population. Retaining gross adjusted cost 6480 and an allowance 1020
makes the next sale easier to follow.

Sell 26 of these units for 2900. IAS~2.34 expenses their current carrying amount,
$26\times 70=1820$. The sale removes 2160 of gross cost and 340 of allowance.
The 52 units still held therefore have gross cost 4320, allowance 680 and carrying
amount 3640. If later evidence of changed circumstances raises net realisable
value to 90 per unit, IAS~2.33 permits a recovery only up to original cost.
Inventory returns to 4320, so the reversal is 680. Reversing the original 1020
would restore 340 relating to goods already sold and would overstate both
remaining inventory and profit.

For the general calculation within this batch, let $q $ be units still held and
$A$ their current allowance. Gross cost is $qc $, and carrying amount is $qc-A$.
A sale of $x $ units removes gross cost $xc $ and allowance $xA/q $, giving
\[
 \text{cost of sales}=x\left(c-\frac{A}{q}\right),
 \qquad A_{\mathrm{remaining}}=\frac{q-x}{q}A.
\]
The first expression charges the carrying amount of the sold goods; the
second keeps only the allowance attributable to unsold goods. At the next
supported valuation $v'$, the required allowance is
\[
 A'=(q-x)\max(c-v',0),\qquad
 \Delta\text{inventory expense}=A'-A_{\mathrm{remaining}}.
\]
A positive difference is a new write-down. A negative difference reduces
inventory expense, subject to the recovery conditions in IAS~2.33. Because
$A'$ cannot be negative, the reduction cannot exceed the allowance for goods
still held. Neither entry affects cash or recreates the cash-flow hedge reserve.

If the final 52 units are sold for 5800 after the recovery, total cost of sales
is $4320+1820+4320=10460$. Net write-down expense is $1020-680=340$, so total
inventory expense remains 10800. Total receipts 14500 give cumulative profit
$90+14500-10800=3790$ in the positive-hedge case and
$-100+14500-10800=3600$ in the negative case. These agree with the earlier
complete-sale results, but the intervening inventory and earnings differ.
Final agreement alone cannot establish that those earlier amounts were correct.

The documentary \code{HedgeInventorySubsequent/1} report reconstructs the
purchase from the retained hedge history, then carries quantity, gross cost
and allowance through dated sales and valuations. Missing recognition or
recovery findings remain unresolved; excess sales, inconsistent dates and
substituted opening amounts are refused. The calculations are cumulative and
pre-tax. They require a homogeneous batch and supplied valuation evidence;
partial forecast purchases, additional inventory lots, tax and company-wide
presentation retain their own methods.\citep[IAS~2, paragraphs 25, 27--34]{eu20231803}
